\chapter{Direct LLVM SSA Compilation \& Native Machine Code (\texttt{shae llvm})}

\section{What is LLVM?}

\begin{kidbox}{The Master Sculptor Robot}
Imagine you draw a blueprint of a bicycle on paper using colored crayons.
Now imagine you hand that paper to the smartest robotic sculptor in the universe.
The robot has diamond laser cutters, supercomputers, and high-tech titanium presses.
The robot studies your blueprint, figures out how to remove extra unneeded metal, calculates the smoothest aerodynamics, and builds an ultra-lightweight racing bicycle that can zoom faster than anything on Earth!
That master robot is called \textbf{LLVM}!
Shae speaks directly to LLVM to turn your friendly code into blazing-fast native machine code!
\end{kidbox}

\section{Static Single Assignment (SSA) Form}

\begin{kidbox}{The Notebook Without an Eraser}
Imagine you are doing math in school, but the teacher takes away your rubber eraser!
Every time you compute a number, you must write it on a brand-new clean line:
\begin{itemize}
    \item Line 1: Let box 1 = 5
    \item Line 2: Let box 2 = 10
    \item Line 3: Let box 3 = box 1 + box 2 (which is 15!)
\end{itemize}
You never cross out or erase an old line! Every box is written once and only once!
This makes it super easy for the computer to optimize, because it never gets confused about which value was in a box at any moment in time!
\end{kidbox}

\begin{deepbox}{LLVM Intermediate Representation (IR) Architecture}
Shae's direct LLVM backend (\texttt{src/llvm.rs} and \texttt{lib/compiler/llvm.shae}) compiles AST representations directly into textual LLVM 18 SSA Intermediate Representation:

\begin{lstlisting}[language=shaelang]
; ModuleID = 'shae_module'
source_filename = "program.shae"
target datalayout = "e-m:e-p270:32:32-p271:32:32-p272:64:64-i64:64-i128:128-f80:128-n8:16:32:64-S128"
target triple = "x86_64-unknown-linux-gnu"

define i64 @add_numbers(i64 %a, i64 %b) {
entry:
  %0 = add nsw i64 %a, %b
  ret i64 %0
}

define i32 @main() {
entry:
  %res = call i64 @add_numbers(i64 100, i64 50)
  call void @print_i64(i64 %res)
  ret i32 0
}
\end{lstlisting}

LLVM performs aggressive optimization passes over this SSA representation:
\begin{enumerate}
    \item \textbf{Constant Propagation \& Folding}: Resolving arithmetic known at compile time into single constants.
    \item \textbf{Dead Code Elimination (DCE)}: Removing unused basic blocks and instructions.
    \item \textbf{Loop Invariant Code Motion (LICM)}: Hoisting computations out of loops.
    \item \textbf{Auto-Vectorization}: Emitting AVX2, AVX-512, and ARM NEON SIMD vector instructions.
\end{enumerate}
Shae invokes Clang with \texttt{-O2} to link the resulting object file directly with the C runtime, producing standalone, zero-dependency ELF and PE executables.
\end{deepbox}

\section{Compiling with \texttt{shae llvm}}

To compile any Shae file to a native standalone executable:
\begin{lstlisting}[language=bash]
shae llvm program.shae -o my_binary
./my_binary
\end{lstlisting}
